\chapter{Problem Statement, Theoretical Basis, and Reference Methods}
\label{ch:introduction}

\section{Problem statement}
Phase-field fracture regularizes a sharp crack over a finite length scale and thereby avoids explicit tracking of a discontinuity. The numerical advantage is accompanied by a spatial-resolution requirement: the diffuse crack zone must be resolved by the finite-element mesh. A uniformly fine mesh is therefore expensive when only a small part of the structure develops damage. The thesis addresses whether Abaqus built-in mesh-refinement capabilities can be combined with a phase-field UEL/UMAT formulation so that fine resolution is concentrated in mechanically relevant regions without changing the underlying fracture model.

The work is organised around four interfaces that are kept explicit throughout the report:
\begin{enumerate}
  \item \textbf{Phase-field constitutive model:} equations and state variables implemented by the supplied user subroutines.
  \item \textbf{Abaqus error estimation and mesh generation:} the standard-element/facsimile representation used to expose stresses and \texttt{MISESERI} to Abaqus.
  \item \textbf{UEL reconstruction:} conversion of the refined physical mesh into the coincident user-element and visualization layers required by the phase-field implementation.
  \item \textbf{State transfer and restart:} transfer of displacement, phase field and history variables when the discretization changes during a continued analysis.
\end{enumerate}
This separation is important because a correct result at one interface does not imply correctness of the complete workflow.

\section{Phase-field fracture formulation}
The implementation follows the standard AT2-type regularization used in the reference phase-field literature \citep{Msekh2015,Molnar2017,Pandey2025}. For displacement field $\mathbf{u}$ and phase field $d\in[0,1]$, with $d=0$ intact and $d=1$ fully broken, the total potential can be written in the generic form
\begin{equation}
\Pi(\mathbf{u},d)=\int_\Omega \left[g(d)\,\psi_0^+(\boldsymbol\varepsilon)+\psi_0^-(\boldsymbol\varepsilon)\right]\,\mathrm{d}\Omega
+\int_\Omega G_c\,\gamma(d,\nabla d)\,\mathrm{d}\Omega-W_{\mathrm{ext}},
\end{equation}
where $G_c$ is the fracture energy, $g(d)$ is the degradation function, and $\gamma$ is the regularized crack-density functional. For the AT2 choice used in the reference work,
\begin{equation}
\gamma(d,\nabla d)=\frac{1}{2}\left(\frac{d^2}{l_0}+l_0\lvert\nabla d\rvert^2\right),
\end{equation}
with regularization length $l_0$. A small residual stiffness is retained through a modified quadratic degradation function, typically $g(d)=(1-d)^2+k$ with $k\ll1$.

Crack irreversibility is enforced through a history field
\begin{equation}
\mathcal{H}(\mathbf{x},t)=\max_{\tau\le t}\psi_0^+(\boldsymbol\varepsilon(\mathbf{x},\tau)),
\end{equation}
so that the crack-driving quantity is non-decreasing. In this thesis, monotonicity of the stored history variable and monotonicity of projected visualization variables are evaluated separately. A decrease in a projected state variable is not automatically interpreted as constitutive healing.

\section{Abaqus UEL/UMAT architecture}
The supplied implementation follows the staggered Abaqus strategy described by Moln\'ar and Gravouil \citep{Molnar2017}. The physical discretization is represented by coincident logical layers. User elements solve the phase-field and mechanical problems, while companion standard elements driven through a UMAT expose stresses and solution-dependent variables to the Abaqus output database. The mapping between these layers is part of the numerical method and is checked explicitly after every regenerated mesh.

The main data path relevant to refinement is
\begin{center}
physical deformation $\rightarrow$ mechanical UEL $\rightarrow$ shared state $\rightarrow$ UMAT/facsimile layer $\rightarrow$ Abaqus stress recovery $\rightarrow$ \texttt{MISESERI}.
\end{center}
Thus \texttt{MISESERI} is a recovered-stress discretization indicator. It is not a phase-field error estimator and is used here only to identify where the coarse mechanical solution has strong unresolved gradients.

The coincident-layer architecture is verified in the input-deck connectivity and element-type audits. A CAE layer excerpt is not inserted here because overlapping layers are not visually separable in the available viewport export; the report does not substitute an illustrative rendering for that model evidence.

\section{Reference mesh-refinement method}
Pandey and Kumar proposed a Python-controlled two-job procedure for phase-field calculations in Abaqus \citep{Pandey2025}. Their script defines a \texttt{RemeshingRule} on the facsimile set \texttt{All\_elem}, requests \texttt{MISESERI}, and subsequently calls \texttt{adaptiveRemesh(odb=...)} using the pre-analysis output database. For the Mode-I edge-crack example, they report a global coarse size of $0.02\,\mathrm{mm}$, local refined size $0.001\,\mathrm{mm}$, $l_0=0.0075\,\mathrm{mm}$, $G_c=2.7\times10^{-3}\,\mathrm{kN/mm}$, and $13{,}941$ elements for the proposed mesh. Their Listing~1 shows \texttt{errorTarget=1.0}, \texttt{refinementFactor=10}, and no coarsening.

The method is therefore described in this report as \emph{Abaqus-native error-guided pre-refinement}. The final fracture calculation uses a fixed refined mesh. This is distinct from automated external-driver remeshing during crack propagation, and from true in-analysis remeshing.

\section{Objectives of the present work}
The work reported here addresses the following questions:
\begin{enumerate}
  \item Can the supplied phase-field implementation reproduce a reference fixed-mesh Mode-I response?
  \item Can the Pandey--Kumar Abaqus-native refinement procedure be reproduced in a headless Python workflow and converted reliably to the layered UEL/UMAT formulation?
  \item Does the refined model reproduce the published force-displacement response and mesh density?
  \item What information must be transferred when the mesh changes during an ongoing phase-field calculation, and does the restarted model preserve mechanical continuity?
  \item Which aspects of the current UEL/UEXTERNALDB architecture limit parallel execution and future ABAQUSER integration?
\end{enumerate}
